% s-018-p018-m002; original PDF 18, 19, 20
\Aphorism{s-018-p018-m002}{9}
\Sperr{Historische Wahrheit.} Das gemalte Portrait ist verdächtig, schon weil es von Zweien erzählt; es ist Gegenstand und Auffassung, Text plus Interpretation. Die Photographie beseitigt wenigstens die Zwischenstation eines Subjects, verlegt aber die Fälschung ins Object zurück: das Gesicht, das stillehalten und eine Platte bestrahlen soll, wird Maske, nämlich starr und unwahr zugleich. Bleibt die Momentaufnahme, die dem Opfer keine Zeit lässt zu posiren und einen „Charakter“ vorzuhängen, die nicht verschmäht indiscret zu sein und Menschlein in ihren fatalsten Augenblicken zu überraschen. Diese Augenblicke auszuwählen ist eine Kunst, die sich wieder mit der des Malers berührt, nur dass dieser sämtliche Momentbilder zu einem einzigen mischt und übereinanderlegt. Aber gerade auf die einzelnen kommt es an, auf die Theile, nicht aufs geistige Band. Wenn die Chemie sich fernerhin um die Ausbildung der Momentphotographie bemüht, kann sie als „Scheidekunst“ auch in psychologicis zu Ehren kommen.
\par

Was wissen wir vom Vergangenen? Grosse Menschen und Ereignisse müssen uns „sitzen“ oder wenigstens Minuten lang, Secunden lang stillehalten, bis sie sich abgezeichnet, eine „Wirkung“ gehabt haben: sie werden gemalt oder photographirt, in beiden Fällen als Kunst, Phantasie, Fälschung überliefert. \Foreign{Il n’y a que l’esprit qui sente l’esprit} ist ein leidlich unwahrscheinlicher Satz: jede Anhängerschaft, jeder Bund zwischen grossen Menschen, jedes Herüber und Hinüber zwischen Mensch und Ereigniss widerlegt ihn. Geist haben privilegirt zum Missverstehen jedes Anderen, das auch Geist hat; Production und Reception schliessen sich aus; grosser Mensch sein heisst grosser Fälscher, grosser „Künstler“ sein. Die Geschichte ist in Gefahr, wenn sie grossen Menschen in die Hände fällt. Und was „Wirkungen“ anbelangt, so wissen wir, dass die Ursachen kein besseres Versteck finden können. Um zu wirken, muss ein Ereigniss sich in seiner natürlichen Beweglichkeit beschränken, muss eine Reihe gleichartiger Erschütterungen aussenden, die sich zu einer „Wirkung“ summiren können, muss, wie gesagt, stillehalten. Und dabei legt es sein Gesicht in fremde Falten und zieht Sonntagskleider an, dabei hat es Wünsche und Absichten und arbeitet der künftigen Retouche vor, dabei lügt es und schämt sich nicht, sich selber zu idealisiren \Dots{3}
\par

Aber es könnte Geschichte geben, die weder Poesie noch Pragmatik wäre, Geschichte in Momentaufnahmen: Menschen und Dinge in ihren fatalsten Augenblicken belauscht, da wo sie, vielleicht zu ihrem eigenen Schaden, aufrichtig sind. Wahrnehmungen Danebenstehender, die der grossen „Wirkung“ (mit der das Falschnehmen beginnt) vorangehen; unschuldige impressions, die für die Impresarien der grossen Geister nicht immer gelegen kamen; kurze verrätherische Gespräche, keine officiellen Reden, keine Bergpredigten und sonstigen Publicitäten, sondern heimliche Zweifel, Einwände, Ketzereien „unter uns“, Dinge, die man in Ecken, wo Zwei \Sperr{miteinander flüstern, auf phonographischem Wege auf}genommen haben würde — wenn es damals Phonographen gegeben hätte. Die gab es nicht, und also fehlt uns diese Gattung Historie: ein Grund mehr, sie nachträglich zu erfinden \Dots{3}
\par

